\section{Gaze Contingency Search
Study}\label{gaze-contingency-search-study}

Experimenter Run Sheet

Version: 0.1 \textbar{} September 16 2026 \textbar{} Protocol proposal
for rehearsal and review. The IRB application has not been submitted.
Use with the approved consent and finalized study settings before
participant collection.

Researcher only: Two counterbalanced voice blocks, generic and
self-similar. Both use gaze-contingent warmer/colder guidance. Planned
trial start is the same fixed front direction. Say ``first block'' and
``second block''; do not suggest that either voice should improve
performance.

Session record: Participant ID \_\_\_\_\_\_ Run \_\_\_\_\_\_ Date
\_\_\_\_\_\_ Researcher \_\_\_\_\_\_ Build \_\_\_\_\_\_ Voice order
\_\_\_\_\_\_ Neutral profile \_\_\_\_\_\_ Survey version \_\_\_\_\_\_

\subsection{Before participant
arrives}\label{before-participant-arrives}

\begin{itemize}
\item
  Clear the swivel-chair rotation area; check floor, cables and nearby
  obstacles. Sanitize and charge the HTC Vive Focus Vision and
  controllers.
\item
  Open the study app and selected survey forms. Confirm coded ID, saved
  voice order, neutral voice profile, audio volume and available
  storage. Never reuse another participant's voice enrollment.
\item
  Record the build and schedule. Implemented beta: two practice trials
  and seven measured trials per voice, 14 total. Final study trial count
  remains open.
\item
  Check the fixed-front gate, matched wording and telemetry against the
  release checklist. These are not all implemented in the current beta;
  do not present a rehearsal as the finalized experiment.
\item
  Prepare consent, baseline survey, both block surveys, final interview
  and session/event record. Use one defined TLX collection route; see
  the beta note in Section 7.
\end{itemize}

\subsection{1 Welcome and consent}\label{welcome-and-consent}

Say: ``Thank you for coming. You will stay seated and turn to find
virtual objects by their color and shape. A voice assistant will provide
hints. You may pause, take a break or stop at any time without
penalty.''

Say: ``We will ask you to record a short voice sample to create a
synthetic version of your voice. This uses cloud voice services. We also
record eye-gaze and task data under a study code. Please read the
consent form for what is collected, where it is processed and how it is
handled.''

ACTION: Review the approved information and obtain consent, including
voice processing authorization. Do not promise offline processing, a
retention period or compensation not specified in the approved
materials. Do not run screen or separate microphone recording unless
separately covered by consent.

\subsection{2 Initial questionnaire and
comfort}\label{initial-questionnaire-and-comfort}

Say: ``Please complete the background and current-comfort questions
before we begin. Use only the study code already entered. Do not add
your name or contact details. You may skip optional questions.''

ACTION: Confirm coded ID and baseline symptom record. Apply approved
screening rules. Record missing responses without inventing answers.
Survey platform and exact form must match the session record.

\subsection{3 Fit and calibrate the
headset}\label{fit-and-calibrate-the-headset}

Say: ``Remain seated in this chair. You may turn the chair, your body
and your head to look around, but do not stand or walk. Tell me
immediately if you feel dizzy, nauseated, strained, uncomfortable or
unsteady.''

ACTION: Fit the Focus Vision, adjust lens spacing/fit using the headset
procedure, enable eye tracking, and complete calibration and validation.
Check that front, side and rear objects are visible and selectable.
Record validation results; do not invent a pass threshold.

Ask: ``Can you see clearly? Is the headset comfortable? Do you feel well
enough to continue?''

ACTION: Dismiss any headset eye-tracking notice, release controls, then
deliberately point to and select Start setup. The app has a separate
start gate to prevent an OS-notice click from starting enrollment.

\subsection{4 Record and prepare the
voices}\label{record-and-prepare-the-voices}

Say: ``First, I will explain the task. You will hear a target color and
shape, then search the objects around you. Hold your gaze on the
matching object to select it. The voice may give warmer or colder
guidance. Reading the passage is for voice preparation, not a scored
search trial.''

Say: ``Read the displayed passage in your natural voice after the tone.
Do not say your name or add personal information. When you finish,
select Finish recording.''

ACTION: Verify the participant ID and selected neutral profile. Let the
app play its recording introduction, countdown and tone. Keep the room
quiet and do not speak over the sample. The current default maximum is
40 seconds; Finish recording allows an earlier finish once enabled.

Implemented app speech: ``We will now connect a short sample of your
voice, after the tone please read the script.'' Then ``Processing
voice.'' When ready: ``Voice setup is ready. Let's check the audio.''

ACTION: Keep the headset active during preparation and verify both voice
playback samples before accepting readiness. ``Ready'' means starter
clips are prepared; the remaining library loads in the background. Do
not continue after failed preparation or substitute the generic voice
for a failed self-similar condition.

\subsection{5 Center the participant and
practice}\label{center-the-participant-and-practice}

ACTION: At the 360 search start checkpoint, seat the participant at the
intended center facing the chosen room-front reference. Release the
trigger, then press Trigger / Enter to center and begin. Record the
reference; do not recenter to wherever the participant happens to face
later.

Say: ``Objects will appear on eight areas around you. Find the object
with the announced color and shape. Turn while staying seated, then hold
your gaze on the matching object until it is selected. You do not need
to count rounds.''

Say: ``The next rounds are practice. They do not count toward the
measured task, although the system may still log them. Use this time to
learn how to search and select an object.''

ACTION: The beta provides two practice rounds across both voices, then a
researcher checkpoint. Check selection and understanding. Correct
controls, not search strategy. Record problems and any additional
familiarization; the final practice/repeat criterion remains open.

Ask: ``How do you select an object? What should you do if you need a
break? Do you have any questions?''

\subsection{6 First voice block}\label{first-voice-block}

Researcher settings: Confirm the saved first voice and planned trial
schedule. Keep object layout, selection method, guidance policy and the
finalized prompt wording constant across voices. The current beta still
varies wording by voice; fix this before the voice-only study.

Say: ``We will now begin the first task block. Find each target as
quickly and accurately as you comfortably can. Stay seated and use the
same gaze selection you practiced. Tell me if you need to pause.''

Planned fixed-front instruction: ``Before each new round, return to the
front reference. Wait for the target instruction and for the objects to
appear, then begin searching.''

ACTION: Use that instruction only with the verified fixed-front gate.
The current beta's transition cue follows the view and does not enforce
front alignment; do not assume the instruction alone provides a
standardized start. Record any rehearsal workaround as a deviation.

ACTION: Start only after the participant is ready. Observe safety and
logging without additional hints or encouragement. The search clock
begins after the instruction and object reveal; reorientation and
preparation are outside active search. Record errors, missing audio and
interruptions.

At completion say: ``That block is complete. Stop turning and relax.
Please answer the next questions about only the block you just
completed.''

\subsection{7 First block questionnaires and
break}\label{first-block-questionnaires-and-break}

Say: ``These questions ask about your workload, your experience of the
activity and the voice assistant. There are no right answers. Please
answer based on the block you just finished.''

ACTION: Collect NASA-TLX, the selected IMI subscales, selected Guo
agent-perception items, voice similarity/ratings and technical checks
using the frozen form. Link every response to participant, run, block,
voice and survey version. Do not administer the full source banks merely
because they are available.

Beta integration note: The app currently requires its in-headset
NASA-TLX submission before the break checkpoint. For a beta rehearsal,
collect this actual TLX response, then remove the headset and administer
the other selected questions. Do not enter dummy TLX values to unlock
the app or collect a second TLX as though it were the same record. The
proposed all-computer survey handoff still requires implementation and
validation.

ACTION: Confirm the questionnaire saved, then hold the researcher
checkpoint while the participant rests. Record break start/end and
comfort. Do not let preparation or a held trigger advance the next
block.

Say: ``You may take a break with the headset off. Tell me when you feel
ready. You may also stop here.''

Before continuing ask: ``Do you feel comfortable and well enough to
continue?''

\subsection{8 Revalidate and run the second voice
block}\label{revalidate-and-run-the-second-voice-block}

ACTION: Refit the headset and revalidate eye tracking and the same
room-front alignment. Record any tracking loss or origin shift. Do not
silently overwrite the origin to correct a mismatch. Confirm the
assigned second voice and remaining audio readiness before releasing the
checkpoint.

Say: ``We will now begin the second task block. The task and selection
method are the same. Find each target as quickly and accurately as you
comfortably can, and tell me if you need to pause.''

ACTION: Use the same trial-start procedure, guidance rules and operator
behavior as in the first block. Verify the second voice actually plays.
Record any technical interruption rather than treating it as poor
participant performance.

At completion say: ``The search task is complete. Please stop turning.
We will finish the questions about this block, then I will help you
remove the headset.''

ACTION: Collect the second block's TLX and other selected measures with
the same route and ordering as the first. Follow the beta TLX gate when
rehearsing that build. Remove the headset safely; record final symptoms
after exposure.

\subsection{9 Final comparison and
interview}\label{final-comparison-and-interview}

Say: ``Please compare your experiences with the two voices. It is fine
to have no preference. We are interested in what you experienced,
including anything distracting, uncomfortable or unhelpful.''

Ask: ``Which voice, if either, did you prefer, and why? What helped or
interrupted your search? Did either voice sound like you? Was anything
unusual or uncomfortable?''

Exploratory follow-up: After open responses, ask ``Did either voice feel
like your own thoughts, like another person helping, or neither?''
Record this as an exploratory report, not proof of an internal-dialogue
mechanism.

ACTION: Use the finalized interview guide and record responses through
the consented method. Confirm final symptom questions and optional
comments. IPQ is on hold; do not automatically administer the old
presence form.

\subsection{10 Debrief and finish}\label{debrief-and-finish}

Say: ``This study compares generic and self-similar voices during
gaze-guided object search. Both voices provide gaze-contingent
assistance. We are examining task performance and how participants
experience the activity and assistant.''

Ask: ``Do you have any questions? Do you feel comfortable and well
enough to leave?''

ACTION: Provide the approved post-participation information, contact
details and any approved compensation. Follow the approved voice-data
handling process; do not promise remote deletion that has not been
verified. Thank the participant and sanitize equipment.

\subsection{Immediate stop or pause
criteria}\label{immediate-stop-or-pause-criteria}

PAUSE OR STOP: Participant request, dizziness, nausea, headache, eye
strain, fatigue, pain, distress, unsafe movement, tracking loss, wrong
voice, inaudible instructions or equipment failure. Help the participant
remain seated and remove the headset when needed.

ACTION: Pause through the verified researcher control; the beta supports
P on a connected keyboard and Trigger / Enter at its resume checkpoint.
App focus loss can pause or technically stop a session. Log time, trial,
reason and outcome. Resume only when the participant wants to continue
and tracking/audio/alignment are valid; otherwise stop and preserve
partial data.

\subsection{Final researcher check}\label{final-researcher-check}

\begin{itemize}
\item
  Confirm both block surveys and final interview/symptom records, or
  record why incomplete. Match participant, run, block and voice
  identifiers.
\item
  Verify gaze\_log.csv, trial\_events.csv, object manifest, trial
  summary, block TLX and voice-readiness/\allowbreak assignment records exist and
  are readable. Inspect actual outcomes; do not code stopped trials as
  completed.
\item
  Secure data using the approved storage/retention workflow. Log
  deviations and technical failures. Keep practice separate from
  measured trials.
\end{itemize}

\subsection{Appendix Current recording
passage}\label{appendix-current-recording-passage}

Implemented in VoiceModeSelector.cs: Read the displayed version, rather
than a different legacy IRB passage. The current text below still
includes shelf and directional language. A natural voice-sample
replacement is planned in issue 25; this is not the finalized study
passage.

``Locate the red sphere. Look toward the upper shelf. Move your gaze
slowly to the right. Check the blue cube near the center. Compare each
object's color and shape. Scan the lower row from left to right. Look
beside the purple cylinder. Focus on the target and hold your gaze
steady. Select the matching object. Take a short pause. Get ready for
the next round.''

\subsection{Appendix Release decisions and source
map}\label{appendix-release-decisions-and-source-map}

Open decisions: Final trials/repetitions, primary outcome, practice
criterion, calibration thresholds, timeout, break duration, final survey
items/scoring and data-retention details. Resolve and record these
before participant collection. No fixed session duration or trial count
is inferred from the meeting.

Implementation gates: Fixed-front start (\#7/\#10), matched voice
wording (\#24), stale-guidance interruption (\#23), validated gaze and
contact timing (\#14/\#20), playback/motion joins (\#21/\#22/\#28),
external survey handoff (\#26/\#27), and reconciled cloud-voice
materials (\#18).

Survey references: NASA-TLX (Hart \& Staveland, 1988); IMI official
complete packet, with Kao et al.~(2021), doi:10.1145/3474665, as
selection precedent; Guo et al.~(2024), doi:10.1145/3651288, Table A1.
Preserve Guo's original attributions to Biocca et al.~(2001), Moussawi
and Koufaris (2019), Zibrek et al.~(2018), Reysen (2005), and Lam et
al.~(2023). SSQ: Kennedy et al.~(1993).

Survey selection: IMI and Guo's survey are included in planning. Full
proposed IMI + Guo ratings + TLX would total 72 ratings per block before
other checks. Appearance-specific Guo questions and overlapping voice
ratings require review. Use docs/measures-register.md for item roles and
scoring decisions.

Project evidence: Assets/VoiceModeSelector.cs and VoiceEnrollment.cs
(capture/acceptance); VoiceSynthesizer.cs (announcements/readiness);
VoicePromptText.cs (wording); SessionConfig.cs (identity/order);
FindObjectGameManager.cs and StudyCheckpoint.cs
(practice/blocks/pauses); TrialDataLogger.cs and GazeDataLogger.cs
(records).

Current documents: docs/STUDY-HUB.md; docs/thesis-study-alignment.md;
docs/measures-register.md; resources/surveys/guo-2024-survey.md. The
older docs/irb/source/voice-script.md still describes offline processing
and must not be used as the current cloud recording script.

Delivery board: https://github.com/users/PedroPovedaQ/projects/2
